\begin{proof}
\textbf{Preliminary cases.}  

If \(n=1\) then both sides of (C) equal \(a_{1}b_{1}\), so the inequality holds with
equality.  

If \(a_{1}=a_{2}= \dots =a_{n}=a\) (or similarly for the \(b\)-sequence), then  

\[
\sum_{i=1}^{n}a_{i}b_{i}=a\sum_{i=1}^{n}b_{i},
\qquad
n\sum_{i=1}^{n}a_{i}b_{i}
   =\Bigl(\sum_{i=1}^{n}a_{i}\Bigr)\Bigl(\sum_{i=1}^{n}b_{i}\Bigr),
\]

so (C) is an equality.  Hence we may assume \(n\ge 2\) and that neither
sequence is constant.

\medskip\noindent\textbf{Reformulation in symmetric form.}  

Define the arithmetic means  

\[
\bar a:=\frac1n\sum_{i=1}^{n}a_{i},
\qquad
\bar b:=\frac1n\sum_{i=1}^{n}b_{i},
\]

and consider  

\[
\Sigma:=\sum_{i=1}^{n}(a_{i}-\bar a)(b_{i}-\bar b). \tag{2.1}
\]

Expanding and collecting terms gives  

\begin{align*}
\Sigma
&=\sum_{i=1}^{n}a_{i}b_{i}
   -\bar a\sum_{i=1}^{n}b_{i}
   -\bar b\sum_{i=1}^{n}a_{i}
   +n\bar a\bar b \\[2pt]
&=\sum_{i=1}^{n}a_{i}b_{i}
   -\frac1n\Bigl(\sum_{i=1}^{n}a_{i}\Bigr)
           \Bigl(\sum_{i=1}^{n}b_{i}\Bigr). \tag{2.2}
\end{align*}

Thus (C) is equivalent to  

\[
\Sigma\ge 0. \tag{2.3}
\]

\medskip\noindent\textbf{Rearrangement inequality argument.}  

\paragraph{Rearrangement Inequality (standard form).}  
If \(x_{1}\le x_{2}\le\cdots\le x_{n}\) and \(y_{1}\le y_{2}\le\cdots\le y_{n}\),
then for every permutation \(\sigma\in S_{n}\),

\[
\sum_{i=1}^{n}x_{i}y_{i}\;\ge\;\sum_{i=1}^{n}x_{i}y_{\sigma(i)}. \tag{3.1}
\]

The identity permutation yields the maximal possible sum.  Applying (3.1) with
\(x_{i}=a_{i}\) and \(y_{i}=b_{i}\) and averaging over all permutations we obtain

\[
\sum_{i=1}^{n}a_{i}b_{i}
   \;\ge\;
   \frac1{n!}\sum_{\sigma\in S_{n}}\sum_{i=1}^{n}a_{i}b_{\sigma(i)}. \tag{3.2}
\]

For a fixed pair \((i,j)\) the term \(a_{i}b_{j}\) appears exactly \((n-1)!\)
times in the double sum, because we may fix \(\sigma(i)=j\) and permute the
remaining \(n-1\) indices arbitrarily.  Hence

\[
\sum_{\sigma\in S_{n}}\sum_{i=1}^{n}a_{i}b_{\sigma(i)}
   =(n-1)!\sum_{i=1}^{n}\sum_{j=1}^{n}a_{i}b_{j}
   =(n-1)!\Bigl(\sum_{i=1}^{n}a_{i}\Bigr)
            \Bigl(\sum_{j=1}^{n}b_{j}\Bigr).
\]

Dividing by \(n!\) gives  

\[
\frac1{n!}\sum_{\sigma\in S_{n}}\sum_{i=1}^{n}a_{i}b_{\sigma(i)}
   =\frac1{n}\Bigl(\sum_{i=1}^{n}a_{i}\Bigr)
               \Bigl(\sum_{j=1}^{n}b_{j}\Bigr). \tag{3.3}
\]

Insert (3.3) into (3.2) and multiply by \(n\):

\[
n\sum_{i=1}^{n}a_{i}b_{i}
   \;\ge\;
   \Bigl(\sum_{i=1}^{n}a_{i}\Bigr)
   \Bigl(\sum_{j=1}^{n}b_{j}\Bigr),
\]

which is precisely inequality (C).  By (2.2) this also shows \(\Sigma\ge0\).

\medskip\noindent\textbf{Equality conditions.}  

The Rearrangement Inequality is strict unless the two sequences are
similarly ordered **and** at least one of them is constant.  Consequently,
equality in (C) occurs iff one of the sequences \((a_{i})\) or \((b_{i})\) is
constant.  When both sequences contain at least two distinct elements the
inequality is strict.

\medskip\noindent\textbf{Conclusion.}  

We have treated the trivial cases, rewritten Chebyshev’s sum inequality in a
symmetric covariance form, and proved its non‑negativity by averaging over all
permutations (a direct consequence of the Rearrangement Inequality).  This
establishes (C) for all \(n\ge 1\) and characterises the equality cases. \(\square\)
\end{proof}